\documentclass[12pt]{article}
\usepackage[utf8]{inputenc}
\usepackage[spanish,es-noshorthands]{babel}
\usepackage{amsmath,amssymb}
\usepackage[margin=2.5cm]{geometry}

\begin{document}

\begin{center}
    {\Large\textbf{Pauta Pregunta 3}}
\end{center}

\vspace{0.4cm}

\noindent\textbf{Datos:}
\[
    I = 1{,}5\ [\text{A}],\qquad
    l_{ab} = 0{,}05\ [\text{m}],\qquad
    l_{bc} = 0{,}13\ [\text{m}],\qquad
    B = 0{,}3\ [\text{T}]\ \text{a } 30^\circ\ \text{sobre } +x
\]

\noindent Tri\'angulo rect\'angulo en $b$ (origen): $a=(0,\;0{,}05)$, $b=(0,0)$,
$c=(0{,}13,\;0)$. La corriente circula $b\to a\to c\to b$ (sentido \textbf{horario}).
Productos cruz de unitarios: $\hat{\imath}\times\hat{\jmath}=\hat{k}$,\;
$\hat{\jmath}\times\hat{\imath}=-\hat{k}$,\;
$\hat{k}\times\hat{\imath}=\hat{\jmath}$,\;
$\hat{k}\times\hat{\jmath}=-\hat{\imath}$.

%==========================================================
\subsection*{(a) Campo magn\'etico $\vec{B}$}

El campo est\'a en el plano $x$--$y$, a $30^\circ$ sobre el eje $+x$:
\[
    \vec{B} = B\cos30^\circ\,\hat{\imath} + B\sin30^\circ\,\hat{\jmath}
            = 0{,}3\cos30^\circ\,\hat{\imath} + 0{,}3\sin30^\circ\,\hat{\jmath}
\]
\[
    \boxed{\vec{B} = 0{,}260\,\hat{\imath} + 0{,}150\,\hat{\jmath}\ [\text{T}]}
\]

%==========================================================
\subsection*{(b) Fuerza magn\'etica sobre cada segmento}

Para cada segmento $\vec{F} = I\,\vec{L}\times\vec{B}$, con $\vec{L}$ en el sentido de la
corriente. Los vectores de cada tramo son:
\[
    \vec{L}_{ab} = 0{,}05\,\hat{\jmath}
    \quad(b\to a),\qquad
    \vec{L}_{ac} = 0{,}13\,\hat{\imath} - 0{,}05\,\hat{\jmath}
    \quad(a\to c),\qquad
    \vec{L}_{cb} = -0{,}13\,\hat{\imath}
    \quad(c\to b)
\]

\noindent\textbf{Segmento $ab$ ($b\to a$):}
\begin{align*}
    \vec{F}_{ab} &= I\,(0{,}05\,\hat{\jmath})\times(0{,}260\,\hat{\imath}+0{,}150\,\hat{\jmath}) \\
                 &= I\big[\,0{,}05\cdot0{,}260\,(\hat{\jmath}\times\hat{\imath})
                    + 0{,}05\cdot0{,}150\,(\hat{\jmath}\times\hat{\jmath})\,\big] \\
                 &= 1{,}5\big[\,0{,}0130\,(-\hat{k}) + 0\,\big]
\end{align*}
\[
    \boxed{\vec{F}_{ab} = -0{,}0195\,\hat{k}\ [\text{N}]}
\]

\noindent\textbf{Segmento $ac$ ($a\to c$, hipotenusa):}
\begin{align*}
    \vec{F}_{ac} &= I\,(0{,}13\,\hat{\imath}-0{,}05\,\hat{\jmath})\times(0{,}260\,\hat{\imath}+0{,}150\,\hat{\jmath}) \\
                 &= I\big[\,0{,}13\cdot0{,}150\,(\hat{\imath}\times\hat{\jmath})
                    - 0{,}05\cdot0{,}260\,(\hat{\jmath}\times\hat{\imath})\,\big] \\
                 &= 1{,}5\big[\,0{,}0195\,\hat{k} - 0{,}0130\,(-\hat{k})\,\big]
                  = 1{,}5\,(0{,}0325\,\hat{k})
\end{align*}
\[
    \boxed{\vec{F}_{ac} = +0{,}0487\,\hat{k}\ [\text{N}]}
\]

\noindent\textbf{Segmento $bc$ ($c\to b$):}
\begin{align*}
    \vec{F}_{cb} &= I\,(-0{,}13\,\hat{\imath})\times(0{,}260\,\hat{\imath}+0{,}150\,\hat{\jmath}) \\
                 &= I\big[\,-0{,}13\cdot0{,}150\,(\hat{\imath}\times\hat{\jmath})\,\big]
                  = 1{,}5\,(-0{,}0195\,\hat{k})
\end{align*}
\[
    \boxed{\vec{F}_{cb} = -0{,}0293\,\hat{k}\ [\text{N}]}
\]

\noindent\textit{Verificaci\'on:}
$\;\vec{F}_{ab}+\vec{F}_{ac}+\vec{F}_{cb}
= (-0{,}0195+0{,}0487-0{,}0293)\,\hat{k} \approx \vec{0}$ \quad(lazo cerrado en campo uniforme).

%==========================================================
\subsection*{(c) Momento magn\'etico $\vec{\mu}$}

\'Area del tri\'angulo:
\[
    A = \tfrac{1}{2}\,l_{ab}\,l_{bc} = \tfrac{1}{2}(0{,}05)(0{,}13) = 3{,}25\times10^{-3}\ [\text{m}^2]
\]
Como la corriente es \textbf{horaria}, por la regla de la mano derecha la normal apunta
hacia adentro de la p\'agina, $\hat{n}=-\hat{k}$:
\[
    \vec{\mu} = I\,A\,\hat{n} = (1{,}5)(3{,}25\times10^{-3})(-\hat{k})
\]
\[
    \boxed{\vec{\mu} = -4{,}875\times10^{-3}\,\hat{k}\ [\text{A}\cdot\text{m}^2]}
\]

%==========================================================
\subsection*{(d) Torque magn\'etico $\vec{\tau}$}

\begin{align*}
    \vec{\tau} = \vec{\mu}\times\vec{B}
       &= (-4{,}875\times10^{-3}\,\hat{k})\times(0{,}260\,\hat{\imath}+0{,}150\,\hat{\jmath}) \\
       &= -4{,}875\times10^{-3}\big[\,0{,}260\,(\hat{k}\times\hat{\imath})
          + 0{,}150\,(\hat{k}\times\hat{\jmath})\,\big] \\
       &= -4{,}875\times10^{-3}\big[\,0{,}260\,\hat{\jmath} + 0{,}150\,(-\hat{\imath})\,\big]
\end{align*}
\[
    \boxed{\vec{\tau} = 7{,}31\times10^{-4}\,\hat{\imath} - 1{,}267\times10^{-3}\,\hat{\jmath}\ [\text{N}\cdot\text{m}]}
    \qquad |\vec{\tau}| = 1{,}463\times10^{-3}\ [\text{N}\cdot\text{m}]
\]

\end{document}
